% 自 notes/docs/1151/seminar/W09_1109_預習.md（87808f3）由 tools/md_to_tex.py 於 2026-11-09 轉出一次；此後本檔為正本。
\chapter{Ferejohn 與 Pasquino〈Constitutional Courts as Deliberative Institutions〉\pg{21 至 36 頁}}
\label{ch:FP}

\begin{longtable}[]{@{}
  >{\raggedright\arraybackslash}p{(\linewidth - 4\tabcolsep) * \real{0.1200}}
  >{\raggedright\arraybackslash}p{(\linewidth - 4\tabcolsep) * \real{0.7500}}
  >{\raggedright\arraybackslash}p{(\linewidth - 4\tabcolsep) * \real{0.1300}}@{}}
\caption{Ferejohn 與 Pasquino 一文的節次與頁次。}\tabularnewline
\toprule\noalign{}
\begin{minipage}[b]{\linewidth}\raggedright
節次
\end{minipage} & \begin{minipage}[b]{\linewidth}\raggedright
小標
\end{minipage} & \begin{minipage}[b]{\linewidth}\raggedright
頁次
\end{minipage} \\
\midrule\noalign{}
\endfirsthead
\toprule\noalign{}
\begin{minipage}[b]{\linewidth}\raggedright
節次
\end{minipage} & \begin{minipage}[b]{\linewidth}\raggedright
小標
\end{minipage} & \begin{minipage}[b]{\linewidth}\raggedright
頁次
\end{minipage} \\
\midrule\noalign{}
\endhead
\bottomrule\noalign{}
\endlastfoot
導論 & （無小標） & 21--22 \\
1 & Deliberative Expectations & 22--25 \\
2 & Democratic Deliberation & 25--27 \\
3 & Constitutional Courts in Modern Democracy（3.1 Constitutional Adjudication in Post-Authoritarian Regimes；3.2 A Special Case: the French Conseil Constitutionel） & 28--33 \\
4 & Deliberation in Constitutional Courts & 33--36 \\
\end{longtable}

\section{導論：Kelsen 型模式的擴散與合法性負擔\pg{21 至 22 頁}}
\label{sec:FP-1}

\subsection{（一）東歐、南歐與戰後歐洲的憲法法院\pg{21 頁，註 1}}

首句即寫「The collapse of the communist empire coincided with a new wave of constitution making」〔共產帝國的瓦解與新一波制憲同時發生〕。東歐各國的新憲法都設有負責憲法審判的機關，其模式是 Ferejohn\ix{persons}{ferejohn@Ferejohn, John} 與 Pasquino\ix{persons}{pasquino@Pasquino, Pasquale} 在後文界定的 Kelsenian\ix{terms}{06@\zhlidxcat{審查制度之類型}!Kelsen 型模式} type（Kelsen\ix{persons}{kelsen@Kelsen, Hans（凱爾森）} 型）。南歐在約二十年前也有相同的現象，起因是希臘、葡萄牙與西班牙的法西斯威權政體瓦解；第二次世界大戰之後，奧地利、義大利與德國走過類似的歷程。Ferejohn 與 Pasquino 據此認為，下文對後威權\ix{terms}{06@\zhlidxcat{審查制度之類型}!後威權政體}政體憲法審判的討論，與東歐的新發展相關\pg{21 頁}。

註 1 補充希臘的情形：希臘自 19 世紀已有司法審查\ix{terms}{01@\zhlidxcat{司法審查權限之歸屬}!司法審查}，1975 年軍事政權結束時制定的憲法，才在分散審查\ix{terms}{06@\zhlidxcat{審查制度之類型}!分散審查與集中審查}制度之上設置一個特別而專門的最高法庭（Special Supreme Tribunal，註 1 所載條次為第 100 條）。

\subsection{（二）模型的共同前提\pg{21 頁，註 2}}

兩人稱這場制度變革可以用「a simple and more or less unified model」〔一個簡單而大致統一的模型〕理解\pg{21 頁}。其前提有二：歐洲各國都致力維持議會對行政與司法部門的權威；就社會學上的事實而言，各後威權\ix{terms}{06@\zhlidxcat{審查制度之類型}!後威權政體}國家對司法（後法西斯國家）、對議會（法國）或對兩者都有不信任。各國因此選擇由設在司法體系之外的專門法院，審查議會的行為。

註 2 建議把兩件事分開看：議會主權\ix{terms}{06@\zhlidxcat{審查制度之類型}!議會主權體制}理論及其階層式的權力分立\ix{terms}{06@\zhlidxcat{審查制度之類型}!權力分立體制}學說，本來就敵視任何司法審查\ix{terms}{01@\zhlidxcat{司法審查權限之歸屬}!司法審查}；制憲者對在威權政體中受教育的法官的不信任，以及對曾與法西斯或共產主義共謀的議會的不信任，則屬於社會學上的期待。兩者合起來，才說明 Kelsen 型\ix{terms}{06@\zhlidxcat{審查制度之類型}!Kelsen 型模式}模式為何成功。

\subsection{（三）合法性負擔與理由說明\pg{21 至 22 頁，註 3 至註 5}}

Ferejohn\ix{persons}{ferejohn@Ferejohn, John} 與 Pasquino\ix{persons}{pasquino@Pasquino, Pasquale} 指出，這個模式把部分立法性的政策形成交到憲法法院手中，「and not merely negative legislative authority as Kelsen\ix{persons}{kelsen@Kelsen, Hans（凱爾森）} himself was forced to admit」〔而且不只是 Kelsen 自己不得不承認的消極立法\ix{terms}{04@\zhlidxcat{法院之功能}!消極立法}權〕\pg{21 頁，註 3 指向 Kelsen 1931 年的 \en{Wer soll der Hüter der Verfassung sein?}，引義大利文譯本 \en{La giustizia costituzionale}, Milano: Giuffré 1981, p. 260}，這使法院承擔額外的 burden of legitimation\ix{terms}{07@\zhlidxcat{憲法審判之審議}!合法性負擔}〔合法性負擔〕。負擔的一部分可由兩種方式處理：使法官隔絕於政治壓力，並准許法院自訂能產生公正或公正外觀的程序，例如祕密開會、以法院名義發布意見\pg{21 至 22 頁}。但除此之外，憲法法院需要為裁判提出有理由的說明；Ferejohn 與 Pasquino 主張憲法法院面對特殊而嚴格的審議期待（deliberative expectations\ix{terms}{07@\zhlidxcat{憲法審判之審議}!審議的期待}）。

這項期待有一部分已寫入創設法院的規範：法國第五共和規範憲法委員會程序的組織法，要求憲法委員會就其裁判提出理由（註 4 引 1958 年 11 月 7 日第 58-1067 號條例第 20 條\ix{statutes}{03@法國!01@憲法委員會組織法（1958 年第 58-1067 號條例）!020@第 20 條}；義大利部分引 1953 年 3 月 11 日第 87 號法第 18 條\ix{statutes}{04@義大利!01@1953 年第 87 號法!018@第 18 條}，判決須「以義大利人民之名」作成並載明理由）。Rawls\ix{persons}{rawls@Rawls, John（羅爾斯）} 在 \en{Political Liberalism}\footcite[231--236]{rawls1993} 中把法院描述為審議機構的典範，即對於國家的強制性政策，期待並提供 reasons, explanations, and justifications〔理由、說明與正當化〕的論壇\pg{22 頁，註 5 指向該書 231 至 236 頁 ``The Supreme Court as Exemplar of Public Reason''}。

\subsection{（四）本章的比較對象\pg{22 頁}}

本章比較歐洲與美國的憲法法院作為審議論壇，對象是法國、德國、義大利、西班牙與美國的法院。兩人主張，這些法院在各政治體系中的位置很不相同：被不同的政治行動者問不同種類的問題，面對不同的期待、歷史與文化及政治上的限制，所以即使處理相當類似的議題，也會採取不同的審議實踐。Ferejohn\ix{persons}{ferejohn@Ferejohn, John} 與 Pasquino\ix{persons}{pasquino@Pasquino, Pasquale} 同時認為，這五個法院都保有 Rawls\ix{persons}{rawls@Rawls, John（羅爾斯）} 所描述的典範性審議特徵\pg{22 頁}。

\section{審議的期待\pg{22 至 25 頁}}
\label{sec:FP-2}

\subsection{（一）審議的界定\pg{22 至 23 頁}}

Ferejohn\ix{persons}{ferejohn@Ferejohn, John} 與 Pasquino\ix{persons}{pasquino@Pasquino, Pasquale} 從 Aristotle\ix{persons}{aristotle@Aristotle（亞里斯多德）} 的審議概念起論：該概念的重點在於審議的目的是批判性地評估、或許改變目標與偏好\pg{22 至 23 頁}。兩人自己的界定放寬了這項要求：「Rather than asking, like Aristotle, whether minds are changed, we would ask only whether public action is conducted through reasoning」〔我們不像 Aristotle 那樣問心意有沒有改變，只問公共行動是否透過說理而進行〕\pg{23 頁}。審議即參與關於公共行動的說理過程，包括對理由保持開放，也就是在有說服力的理由出現時願意改變偏好、信念或行動，並願意提出理由說服他人，排除威脅強制或欺瞞。

Ferejohn 與 Pasquino 承認美國最高法院的成員很少改變心意，至少在深層的政治承諾上如此。但法院的裁判作用於真實或假設的案件，目的在決定具體情境中可以或應該做什麼，所以對基本原則有歧見的法官，可以在較具體的層次上取得一致：「While abstract ideas may not be open to revision, ideas about particular cases may be much more flexible and pragmatic」〔抽象的觀念或許無從修正，關於個案的想法卻可以靈活得多、務實得多〕\pg{23 頁}。

\subsection{（二）理由說明的兩層意義\pg{23 頁，註 6}}

Ferejohn\ix{persons}{ferejohn@Ferejohn, John} 與 Pasquino\ix{persons}{pasquino@Pasquino, Pasquale} 指出，高等法院不同於其他政治機構，不只公布命令或裁判，也被期待公布其主文的 plausible rationales〔可信的理據〕，即他人不論自己的利益是否得到伸張，都可以期待加以尊重\ix{terms}{05@\zhlidxcat{解釋理論與自制}!尊重}與接受的論證\pg{23 頁}。這項期待在兩個意義上屬於審議的期待\ix{terms}{07@\zhlidxcat{憲法審判之審議}!審議的期待}。

\begin{enumerate}
\def\labelenumi{\arabic{enumi}.}
\tightlist
\item
  法院所給的理由可以被理解並接受為在某種規範意義上「我們自己的」行動理由；用 Rawls\ix{persons}{rawls@Rawls, John（羅爾斯）} 的語彙說，法院提供 public reasons〔公共理由\ix{terms}{07@\zhlidxcat{憲法審判之審議}!公共理由}〕。
\item
  法院是合議制機構，這些理由經由受該法院的裁判規範所引導的內部審議\ix{terms}{07@\zhlidxcat{憲法審判之審議}!內部審議}過程取得。這個過程是否受法院將公布單一意見或同時公布多數意見的共同期待規範，視情形而定。法院所公布的理由，在這個意義上是在一個「ranging from consensus seeking to majoritarian」〔從追求共識到多數決〕的規範框架內協商出來的。
\end{enumerate}

註 6 以美國最高法院愈來愈常公布不同意見\ix{terms}{04@\zhlidxcat{法院之功能}!不同意見與協同意見}書與協同意見書為例，說明這些裁判規範具有歷史上的彈性；註中另提出一種可能的解釋：在日益多元的文化中，對主文提出多種理由說明\ix{terms}{07@\zhlidxcat{憲法審判之審議}!理由說明}或有好的理由。

\subsection{（三）理由說明對法院權威的兩項功能\pg{24 至 25 頁}}

兩人認為理由說明\ix{terms}{07@\zhlidxcat{憲法審判之審議}!理由說明}對法院權威特別重要，分兩方面\pg{24 頁}。第一方面稱為民主的：由於（多數）法官不受選舉問責，裁判中所提出的理由，特別是其中一些理由植根於選舉問責機關的實定行為時，可以為公共行動提供間接的民主正當化。即使司法審議無法追溯到立法行為，只要植根於民主政府所由以建立、或民主的人民所預設的憲法或道德原則，司法理由仍可視為間接或遞移地民主。Ferejohn\ix{persons}{ferejohn@Ferejohn, John} 與 Pasquino\ix{persons}{pasquino@Pasquino, Pasquale} 因此認為，審議的法官可能增強民選官員的權力，替他們提供更有彈性且更聰明的制度，以執行並精練自己的立法；若這是立法者對法官的期待，民主的立法者就應該樂於擴張司法權。Ferejohn 與 Pasquino 另提出相反的可能：依憲法或道德原則行事的審議法官，可能增強人民（作為集體或個人）相對於其民主代表的權威。Rawls\ix{persons}{rawls@Rawls, John（羅爾斯）} 的公共理由\ix{terms}{07@\zhlidxcat{憲法審判之審議}!公共理由}說，在這個意義上強調了法官的民主性格，因為公共理由是不論各人對善的整全觀點為何，都可期待其認同的理由。司法審議對立法者也有風險：法官依理由行事，可以修改或廢棄民主的命令，並要求立法模式具有某種一致性或合理性，禁止民選官員制定某些政策，不論這些政策在政治上多麼有吸引力。這些風險是贊成民主或反民主，兩人認為取決於立法機關與政治共同體之間的連結等因素。

第二方面是司法所提供的理由使國家官員、其他法官、律師與一般公民，得以預期本次裁判對日後案件的意涵\pg{24 頁}。這項功能 ``efficiency enhancing''〔提升效率〕，發揮預測與協調的作用，不具正當化的作用；它使法院得以引導或協調大量私人與公共行為，而不致有太多徒勞的訴訟\pg{24 頁}。Ferejohn 與 Pasquino 接著說，效率也可以構成正當化，但「this kind of justification has nothing to do with democracy and would be available to courts in any legal system」〔這種正當化與民主無關，在任何法律體系中法院都可以援用〕\pg{24 至 25 頁}。

審議與說理是任何法律體系中法官的合法性來源，但在民主政體中分量特別重：政治權威須被視為直接或間接來自人民，民主機構須回應民意，這兩項期待使法官處於「a peculiarly weak position」〔特別脆弱的地位〕，這也是對司法審查\ix{terms}{01@\zhlidxcat{司法審查權限之歸屬}!司法審查}的 countermajoritarian\ix{terms}{05@\zhlidxcat{解釋理論與自制}!反多數決困難}〔反多數決〕面向感到憂慮的根源\pg{25 頁}。

\section{民主制度中的審議\pg{25 至 27 頁}}
\label{sec:FP-3}

\subsection{（一）行政、立法與選舉所受的審議期待\pg{25 至 26 頁，註 7}}

Ferejohn\ix{persons}{ferejohn@Ferejohn, John} 與 Pasquino\ix{persons}{pasquino@Pasquino, Pasquale} 拿法院所受的期待與其他機構對照：「In contrast to what we expect from courts, we do not demand as much by way of reason-giving from our other, more evidently democratic, institutions」〔相對於我們對法院的期待，我們對其他更明顯屬於民主的機構，沒有要求那麼多的理由說明\ix{terms}{07@\zhlidxcat{憲法審判之審議}!理由說明}〕\pg{25 頁}。三類機構的情形如下表。

\begin{longtable}[]{@{}
  >{\raggedright\arraybackslash}p{(\linewidth - 4\tabcolsep) * \real{0.1000}}
  >{\raggedright\arraybackslash}p{(\linewidth - 4\tabcolsep) * \real{0.8300}}
  >{\raggedright\arraybackslash}p{(\linewidth - 4\tabcolsep) * \real{0.0700}}@{}}
\caption{行政、立法與選舉所受的審議期待。}\tabularnewline
\toprule\noalign{}
\begin{minipage}[b]{\linewidth}\raggedright
機構
\end{minipage} & \begin{minipage}[b]{\linewidth}\raggedright
兩人所述的審議期待
\end{minipage} & \begin{minipage}[b]{\linewidth}\raggedright
頁次
\end{minipage} \\
\midrule\noalign{}
\endhead
\bottomrule\noalign{}
\endlastfoot
行政機關 & 可期待未經選舉的官僚在其決定與立法機關的決定之間提供某種理性連結；在美國及多數西方民主國家，不要求行政行為的理由既清楚又有說服力，只要所用手段與法律所定目的大致成比例；對個別主張作判斷時，正當程序的考量使行政機關承受類似法院的審議期待，制定一般或抽象政策時，這項期待改由選舉程序承擔（註 7） & 25 \\
立法機關 & 更少要求立法機關在任何法律論壇中說明理由；政黨綱領遠不足以為具體立法提案提供嚴謹的理由；法律文本的前言很少被認真對待，也幾乎不曾實質改變法律本身的理解；容許立法者互相護航（logroll）、拼湊式立法（checkerboard statutes）、聖誕樹式法案（Christmas tree bills），法律上只限制終局的立法行為與憲法價值或規範相衝突的情形 & 25--26 \\
選舉 & 公共規範近乎禁止說理，例如祕密投票、禁止在投票所附近競選、一人一票的民主觀，以及對競選捐款與支出的限制；法院不願追查選舉結果（含公民創制與公民投票）背後不被允許的意圖 & 26 \\
\end{longtable}

Ferejohn 與 Pasquino 把投票所稱為「a reason free zone\ix{terms}{07@\zhlidxcat{憲法審判之審議}!無須理由的領域}」〔無須理由的領域〕，行使投票權不以選民能提出或擁有好的理由為條件；理由若要反映在選舉程序中，必須透過個別選民，說服她在兩組考量之間擇一投票。Ferejohn 與 Pasquino 的結論是：「it is numbers and not reasons that count」〔計算的只有票數，理由不在其內〕\pg{26 頁}。

\subsection{（二）對公共審議可被腐化的疑慮\pg{26 至 27 頁}}

上述限制的理由，兩人認為相當明顯：自 Plato\ix{persons}{plato@Plato（柏拉圖）} 以來的理論家，擔心公共審議可被腐化，普通公民太容易被誘人的修辭或賄賂左右，或被公職人員與私人利益所脅迫\pg{26 頁}。歷史上這些疑慮引起對政黨、利益團體等可能壓倒普通公民的組織的深刻懷疑。Rousseau\ix{persons}{rousseau@Rousseau, Jean-Jacques（盧梭）} 與 Madison\ix{persons}{madison@Madison, James（麥迪遜）} 是兩個顯著的例子：Rousseau 與 Madison 都認為公共領域可被腐化，因而竭力尋求能引導這些強大傾向的制度改革。Rousseau 在《社會契約論》中主張，若可能，應禁止政黨與其他派系的形成，並禁止立法提案以私人個人與團體為對象；他雖然惋惜羅馬人必須採用祕密投票，仍承認共和國的腐化已太深，不容公開行使投票權\pg{26 至 27 頁}。Madison 對直接限制或規制私人利益持較懷疑的態度，提議以增加利益的數量來限制其在公共論壇的影響力。Ferejohn\ix{persons}{ferejohn@Ferejohn, John} 與 Pasquino\ix{persons}{pasquino@Pasquino, Pasquale} 記載，美國人又過了一百年才採用祕密投票，以限制選舉權的腐化\pg{27 頁}。

\subsection{（三）期待落差的三種解釋\pg{27 頁}}

Ferejohn\ix{persons}{ferejohn@Ferejohn, John} 與 Pasquino\ix{persons}{pasquino@Pasquino, Pasquale} 概述民主制度中各機構所受的審議期待：最民主、即最貼近人民的機構，完全不被期待審議，反而被各種妨礙審議的條件包圍；最不民主的決策機構，則被期待成為近乎純粹的公共理由\ix{terms}{07@\zhlidxcat{憲法審判之審議}!公共理由}論壇\pg{27 頁}。兩人提出三種可能的解釋。

\begin{enumerate}
\def\labelenumi{\arabic{enumi}.}
\tightlist
\item
  審議或說理可視為較直接的民主正當性的替代品：有票數就不必說服他人跟隨，因為他們選你就已承諾跟隨；沒有票數，說服與說理就是唯一的辦法。審議與說理因此是非民主機構取得人們服從的方式，是法院與行政機關為作成權威性決定而必須承擔的特別負擔。
\item
  西方民主傳統深處有一種懷疑，認為審議要求在某種程度上反民主或菁英主義。告訴人民或立法機關只有多數不夠，還必須有好的理由才能行動，等於弱化民主生活的一項要素。這個立場未必反對審議，但對在多數行動上加諸審議要求持懷疑態度；理由若是好的，自會影響多數並反映於票數。Ferejohn 與 Pasquino 寫道：「Reasons, in a democracy, should not get counted twice」〔在民主之下，理由不應被計算兩次〕。
\item
  在憲法民主之下，公共政策的選擇發生在規範的位階之中：道德與憲法規範適當地限制民主的選擇。規範這些規範與民主選擇之間的界線時，說理與審議特別重要，因為所設想的是告訴人民或其代表，某項政策不被允許。法院與行政機關（當其行為如法院時），因此在憲法價值可能涉及的領域承擔審議期待；立法機關在作為憲法解釋者時，會面臨類似的期待，但在制定明確處於憲法可接受範圍之內的政策時則不然。
\end{enumerate}

\section{憲法法院在現代民主中的位置\pg{28 至 33 頁}}
\label{sec:FP-4}

\subsection{（一）議會主權體制的審查\pg{28 頁，註 8、9}}

Ferejohn\ix{persons}{ferejohn@Ferejohn, John} 與 Pasquino\ix{persons}{pasquino@Pasquino, Pasquale} 說：「Roughly speaking, we may describe the institutional position of constitutional courts as situated along a single dimension」〔大致而言，可以把憲法法院的制度位置描述為落在單一向度上〕，左端是純粹的議會主權\ix{terms}{06@\zhlidxcat{審查制度之類型}!議會主權體制}體制，例如英國與法國第三、第四共和，右端是 Montesquieu\ix{persons}{montesquieu@Montesquieu（孟德斯鳩）} 式的權力分立\ix{terms}{06@\zhlidxcat{審查制度之類型}!權力分立體制}體制，中間是混合立法優越與其他較獨立機構的各種憲法體制\pg{28 頁}。議會主權體制依定義把行政與法院置於立法機關之下，其角色在執行並實施立法者的命令；司法審查\ix{terms}{01@\zhlidxcat{司法審查權限之歸屬}!司法審查}的觀念因此與這類法制完全格格不入。

圖~\ref{fig:FP-placement} 依第 3 節並列三種體制中各機關的配置、憲法審查權之所在與審查介入之時點。

\begin{qafigure}[三種體制中憲法審查權之所在]{三種體制中憲法審查權之所在。Ferejohn 與 Pasquino，28 至 33 頁。\label{fig:FP-placement}}
% Ferejohn 與 Pasquino 第 3 節（28 至 32 頁）：三種體制中憲法審查權之所在。
% 欄序依兩人所說的單一向度，左端議會主權、右端權力分立（28 頁）；Kelsen 型居中係依 30 頁「departed from the
% hierarchical parliamentary model」之描述排列。垂直、水平二詞出自 29 頁「separated horizontally rather than vertically」。
% 色調：審查權所在之機關 accent（terracotta），其餘機關 slate；卡片淡底加左側色條（skill diagram-design「動手前」第 2 條）。
% 各格文字逐一回刊本核過頁次（2026-11-09）。
\definecolor{figRev}{HTML}{B08060}\definecolor{figRevbg}{HTML}{F5ECEB}
\definecolor{figLine}{HTML}{E8D3D1}
% 卡片：\placecard{色}{左}{上}{右}{下}{內容}
\providecommand{\placecard}[6]{%
  \fill[#1bg] (#2,#3) rectangle (#4,#5);
  \fill[#1] (#2,#3) rectangle ($(#2,#5)+(0.6,0)$);
  \node[t] at ($(#2,#3)+(2.2,-1.4)$) {#6};}
\begin{tikzpicture}[qa,x=1mm,y=1mm,
    t/.style={anchor=north west,inner sep=0pt,align=left,execute at begin node={\raggedright\setstretch{1.2}}},
    cell/.style={t,font=\scriptsize,text width=40mm},
    rowlbl/.style={t,font=\scriptsize,text=dimink,text width=17mm},
    link/.style={figGray,line width=0.5pt},
    arr2/.style={figGray,line width=0.6pt,-{Stealth[length=4pt,width=3.2pt]}}]
  % 單一向度
  \draw[figGray,line width=0.5pt,{Stealth[length=3.5pt,width=3pt]}-{Stealth[length=3.5pt,width=3pt]}] (19,6) -- (146,6);
  \node[font=\scriptsize,text=dimink,anchor=south] at (82.5,6.6) {制度位置的單一向度（28 頁）};
  % 欄頭
  \foreach \x/\h/\s in {
      19/議會主權體制/{垂直：行政與法院\\從屬於立法機關（28 頁）},
      62.5/{Kelsen 型模式}/{憲法法院置於司法與\\其他部門之外（30 頁）},
      106/權力分立體制/{水平：三權分立（29 頁）}}{
    \node[t,font=\footnotesize] at (\x,2.5) {\h};
    \node[t,font=\scriptsize,text=dimink] at (\x,-1.8) {\s};
  }
  % 列名與列間淡線
  \node[rowlbl] at (0,-13) {機關之配置};
  \node[rowlbl] at (0,-50) {審查之時點\\與方式};
  \node[rowlbl] at (0,-61.5) {違憲之效果};
  \node[rowlbl] at (0,-73) {例};
  \foreach \y in {-10,-47.5,-59,-70.5}{\draw[figLine,line width=0.4pt] (0,\y) -- (146,\y);}
  % 議會主權：審查在立法機關內部，行政與法院居其下
  \begin{scope}[shift={(19,0)}]
    \placecard{figGray}{0}{-13}{40}{-34}{立法機關}
    \placecard{figRev}{3.5}{-19}{38}{-32}{憲法審查在其內部\\{\scriptsize 上議院、司法委員會、}\\{\scriptsize 法國憲法委員會（28 頁）}}
    \placecard{figGray}{0}{-38.5}{18}{-44.5}{行政}
    \placecard{figGray}{22}{-38.5}{40}{-44.5}{法院}
    \draw[link] (9,-34) -- (9,-38.5);
    \draw[link] (31,-34) -- (31,-38.5);
  \end{scope}
  % Kelsen 型：憲法法院居各部門之外、之上；普通法院停止訴訟並移送
  \begin{scope}[shift={(62.5,0)}]
    \placecard{figRev}{0}{-13}{40}{-29}{憲法法院\\{\scriptsize 審查權集中於此，}\\{\scriptsize 居各機關之上（31 頁）}}
    \placecard{figGray}{0}{-38.5}{10}{-44.5}{立法}
    \placecard{figGray}{12}{-38.5}{22}{-44.5}{行政}
    \placecard{figGray}{24}{-38.5}{40}{-44.5}{普通法院}
    \draw[arr2] (32,-38.5) -- (32,-29);
    \node[t,font=\scriptsize,text=dimink,anchor=north east,align=right] at (30.5,-30.6) {疑法律違憲者\\停止訴訟並移送};
  \end{scope}
  % 權力分立：三權並列，審查權散布於各級法院
  \begin{scope}[shift={(106,0)}]
    \placecard{figGray}{0}{-13}{10}{-19}{立法}
    \placecard{figGray}{12}{-13}{22}{-19}{行政}
    \placecard{figGray}{24}{-13}{40}{-19}{司法}
    \placecard{figRev}{24}{-25.5}{40}{-31.5}{最高法院}
    \placecard{figRev}{24}{-38.5}{40}{-44.5}{下級法院}
    \draw[link] (32,-19) -- (32,-25.5);
    \draw[link] (32,-31.5) -- (32,-38.5);
    \node[t,font=\scriptsize,text=dimink,anchor=north east,align=right] at (22.5,-27.6) {任何法院皆得審查，\\最高法院居其首};
  \end{scope}
  % 各列內容
  \foreach \x/\a/\b/\c in {
      19/{事前、抽象：比較法律\\與憲法之文本（28 頁）}/{法國：提案不得公布；\\公布後不受審查（32 頁）}/{英國、芬蘭、瑞典；法國\\第三至第五共和（28 頁）},
      62.5/{抽象審查之外，兼有\\事後之具體審查（31 頁）}/{廢止法律並重建法律狀態，\\兼為消極與積極立法（30 頁）}/{奧地利、德國、義大利、\\西班牙、葡萄牙（30 頁）；\\東歐（31 頁）},
      106/{於具體案件中為之，\\非抽象、非事前（29 頁）}/{法律不廢止，僅於該案\\拒絕適用（29 頁）}/{美國、挪威（29 頁，註 11）}}{
    \node[cell] at (\x,-50) {\a};
    \node[cell] at (\x,-61.5) {\b};
    \node[cell] at (\x,-73) {\c};
  }
\end{tikzpicture}

\end{qafigure}

兩人推論，議會主權體制必須把憲法審查設在立法機關之內，否則議會就不是主權者。因此憲法審查只能是 a priori\ix{terms}{06@\zhlidxcat{審查制度之類型}!抽象審查與具體審查}（先於立法提案公布為法律）而且是 abstract（僅涉及法律與憲法文本的比較）；援引憲法原則的權限，須集中於立法機關內部少數人手中。Ferejohn 與 Pasquino 舉英國、法國、芬蘭與瑞典（直到最近）為例，憲法審查全部發生在立法機關之內。傳統上這項審查可集中於上議院（可能有某種消極的立法權，如停止性否決權）、某種司法委員會，或獨立的機構如法國的國務院（Conseil d'Etat），後者可就憲法問題向立法機關提供意見。法國第五共和更設計出「a third legislative body」〔第三個立法機構〕，即 Conseil Constitutionel，可以修改或否決立法提案（註 8 引 Alec Stone, \en{The Birth of Judicial Politics in France}, Oxford 1992，與 John Bell, \en{French Constitutional Law}, Oxford 1992）。這個機構原由戴高樂派（Gaullists）創設，用來使議會合理化或加以控制，後來演變為允許對政府的立法提案作一般性憲法審查。無論審查或意見來自何處，把憲法原則適用於立法的權限仍在立法機關自身\pg{28 頁}。

註 9 記載這項轉變的兩個關鍵時刻：1971 年憲法委員會主張有權依廣泛而未成文的憲法原則（所謂 bloc de constitutionnalité\ix{terms}{06@\zhlidxcat{審查制度之類型}!憲法整體}〔憲法整體〕）廢棄政府的立法，在該案中從「共和國基本原則」辨識出結社自由的權利，並把 1789 年人權宣言與 1946 年憲法前言提升為可用以推翻立法提案的憲法規範；1974 年，政府擔心選舉失敗，成功促成一項修憲\ix{terms}{03@\zhlidxcat{憲法之效力基礎}!修憲}，使任何 60 名議員得將法律提交憲法委員會，此後憲法委員會憑 1971 年創造的強力憲法工具，愈來愈集中審查\ix{terms}{06@\zhlidxcat{審查制度之類型}!分散審查與集中審查}政府的提案\pg{28 至 29 頁}。

\subsection{（二）分散而具體的司法審查\pg{29 至 30 頁，註 10、11}}

權力分立\ix{terms}{06@\zhlidxcat{審查制度之類型}!權力分立體制}體制依定義把立法、行政、司法三權視為水平分立，並傾向把憲法權限賦予行使各項權力的機構\pg{29 頁}。Ferejohn\ix{persons}{ferejohn@Ferejohn, John} 與 Pasquino\ix{persons}{pasquino@Pasquino, Pasquale} 以 Montesquieu\ix{persons}{montesquieu@Montesquieu（孟德斯鳩）} 為例：他把無權力分立的狀態定義為專制，並認為法院適用法律於個案時，不應有任何重要意義上的立法作用，司法權因此是 pouvoir nul。兩人補充，這不表示法官（或更精確地說，陪審團）在適用法律時全無裁量；法院可以因適用某法律會產生不正義而拒絕在特定案件中適用，該法律仍然存在，其效力僅在該案中被排除。

這種解釋法律與事實以便適用的權力，產生一種特色鮮明的司法審查\ix{terms}{01@\zhlidxcat{司法審查權限之歸屬}!司法審查}，與議會主權\ix{terms}{06@\zhlidxcat{審查制度之類型}!議會主權體制}體制不同\pg{29 頁}。Ferejohn 與 Pasquino 列出三點：審查權散布於整個司法體系，因為任何法院為適用法律都須解釋憲法或制定法；審查在具體案件中行使，既非抽象也非 a priori\ix{terms}{06@\zhlidxcat{審查制度之類型}!抽象審查與具體審查}；審查不具立法性，法律並未被廢止，只是對該案（視法制而定，也及於類似案件）拒絕適用。Ferejohn 與 Pasquino 以美國憲法第三條為例，該條使最高法院居於聯邦司法部門之首，只准許憲法審查產生於真正的案件或爭議，並准許任何法院進行這類審查\pg{29 頁}。兩人寫道，在這方面美國的憲法審判制度與法國及多數歐洲模式「completely different」〔完全不同〕；法院適用制定法時必須常以憲法讀之，絕不以違反憲法保障的方式適用，所以制定法常被賦予使其配合憲法架構的解釋，偶爾被完全拒絕適用。這些作為也許有類似立法的效果，但不像法國制度那樣在形式上改變或消除制定法文本；違憲的制定法並未被廢止，只是不適用於特定爭議。Ferejohn 與 Pasquino 說明，這樣的說法凸顯司法審查的司法面向\pg{29 至 30 頁}。

註 10 指出，司法審查的非立法性質是理論上的觀念；實行分散審查\ix{terms}{06@\zhlidxcat{審查制度之類型}!分散審查與集中審查}的司法體系，通常會設計對下級法院的階層控制，從歧異的裁判中形成一致的規則或法理，這使司法行為更具立法性。註 11 提到挪威的憲法審判制度在許多方面與美國相似：審查發生於具體案件，且分散於普通法院\pg{29 頁}。

\subsection{（三）Kelsen 型模式的立法功能論\pg{30 至 31 頁，註 12、13}}

Ferejohn\ix{persons}{ferejohn@Ferejohn, John} 與 Pasquino\ix{persons}{pasquino@Pasquino, Pasquale} 寫道，歐洲多數法律體系有議會主權\ix{terms}{06@\zhlidxcat{審查制度之類型}!議會主權體制}的傳統，即行政與司法皆隸屬於議會，但 20 世紀的發展，尤其兩次世界大戰之後，侵蝕了這些承諾；因此一次大戰後的奧地利，二次大戰後的德國與義大利（以及法西斯政權瓦解後的西班牙與葡萄牙），採用了在重要方面偏離階層式議會模型的憲法\pg{30 頁}。兩人為方便起見把這些體制統稱 Kelsenian\ix{terms}{06@\zhlidxcat{審查制度之類型}!Kelsen 型模式}，取名自發明其獨特制度形式的奧地利法學家。

Ferejohn 與 Pasquino 歸納 Kelsen\ix{persons}{kelsen@Kelsen, Hans（凱爾森）} 的核心觀念為：憲法審判與其說是司法功能，不如說是立法功能。憲法法院宣告制定法違憲，不只是在消極意義上立法而廢止一部法律，由於必須重建該制定法之前的法律狀態，也在積極意義上立法\pg{30 頁，註 13}。註 13 說明，憲法法院依憲法價值解釋制定法時更是如此；Kelsen 本人對憲法法院角色的界定曾經猶豫，起初稱為「negative legislation」〔消極立法\ix{terms}{04@\zhlidxcat{法院之功能}!消極立法}〕，後來回應 C. Schmitt\ix{persons}{schmitt@Schmitt, Carl（施密特）} 時，在 \en{Wer soll der Hüter der Verfassung sein?}（1931）中承認法院有積極的立法功能（Ferejohn 與 Pasquino 另引 P. Pasquino, ``Gardien de la constitution ou justice constitutionnelle? C. Schmitt et H. Kelsen'', in \en{1789 et l'invention de la constitution}, Paris: Bruylant, 1994, pp. 141--152）。註 12 記載奧地利於 1920 年成為歐洲大陸第一個設憲法法院的國家，Hans Kelsen 在其設立中扮演關鍵角色，其歷史前身是奧匈帝國與神聖日耳曼帝國（Reichskammergericht）中裁判中央政府與各邦（Länder）衝突的法院。

依兩人的描述，Kelsen 對憲法審判的理解著重其立法面向，把它設想為制定法與憲法文本的比較；這種 abstract review 產生於文本的 a priori\ix{terms}{06@\zhlidxcat{審查制度之類型}!抽象審查與具體審查} 抽象比較，與有真實（受害）當事人主張權利的具體個案無關\pg{30 頁}。這樣看待的憲法審判「seems inherently political」〔似乎本質上具有政治性〕，因為憲法法院必須在規範社會行為的各種替代規範規則之間審議並選擇。Kelsen 因此認為憲法法院應置於司法以及其他政府部門之外，其權力由經政治任命的法官行使，這些法官通常來自特別能作文本抽象比較、並有能力就規範審議且說明裁判的人，未必來自有司法經驗者\pg{30 至 31 頁}。

\subsection{（四）後威權政體的憲法審判\pg{31 至 32 頁，註 14 至 16}}

在 3.1 中，Ferejohn\ix{persons}{ferejohn@Ferejohn, John} 與 Pasquino\ix{persons}{pasquino@Pasquino, Pasquale} 寫道 Kelsen\ix{persons}{kelsen@Kelsen, Hans（凱爾森）} 的觀念對後威權\ix{terms}{06@\zhlidxcat{審查制度之類型}!後威權政體}政體特別有吸引力：奧地利、德國、義大利與西班牙都採行，1989 年之後也在東歐生根\pg{31 頁}。各後威權憲法所設的憲法審判機構，可審查立法行為，有時並審查行政與司法行為。但各國都正處於從舊的、被不信任的政權過渡到新政權的過程，法官因而被特別懷疑，不宜作為憲法審查權的持有者（註 14 說明這些國家的司法體系基本上是封閉的職業階層，特別不受外界影響，這種極端的隔離使它們在威權崩潰之後更不可能被賦予重要的額外權力）。所以有強有力的政治理由，把憲法審判置於司法之外：審查機構實際上位於其他各機構之上，得從憲法觀點審查任何政府行為，並交給經政治任命的專門機構。

Ferejohn 與 Pasquino 指出，Kelsen 強調抽象的憲法審查，但所有現代後威權憲法法院也都有具體的 a posteriori 審查權\pg{31 頁}。這些法院的管道，不只由政府與政治少數控制，也由個案中的普通訴訟當事人，或如義大利由普通法院控制。因此法院可能被要求以 Kelsen 式的方式抽象比較憲法與制定法文本（由直接提交憲法問題），或面對下級法院在進行中的案件所生的憲法問題，或面對整個已裁判的案件（德國與西班牙均有，註 15 說明對設於 Strasbourg 的歐洲人權法院的申訴，傳統上也如此，須待訴訟當事人用盡國內救濟後才審查）。在任何這些情形，廢棄或修改立法的實際權限，一般集中於憲法法院，不分散於整個司法體系。普通法院若懷疑法律合憲性，必須停止其面前的程序，並將問題提交憲法法院決定\pg{31 頁；註 16 說明適用歐洲共同體法的法院遵循相同的做法}。

憲法法院的組成也特別\pg{31 至 32 頁}。權力分立\ix{terms}{06@\zhlidxcat{審查制度之類型}!權力分立體制}體制傾向分散而具體的審查，普通法官可望發展出憲法審判的能力；在審查集中而且常為抽象的後威權體制，普通法官並無特殊的權威主張，加上威權的過去，法官至少在初期被懷疑能否擔任憲法與民主價值的仲裁者。因此在所有後威權體制中，法學教授往往與若干法官共同占據法院的多數席位\pg{32 頁}。

\subsection{（五）法國憲法委員會的特殊位置\pg{32 至 33 頁，註 17}}

3.2 把法國第五共和的憲法審查獨立討論。兩人說法國的共和傳統穩固屬於議會主權\ix{terms}{06@\zhlidxcat{審查制度之類型}!議會主權體制}傳統，自大革命以來即敵視憲法審判；但戴高樂及其支持者堅持對議會設置制度性的限制，憲法委員會即其一\pg{32 頁；註 17 記載 1795 年後雅各賓的國民公會一致否決了 Sieyès 提出的 jury constitutionnaire〔憲法陪審〕提案}。在議會主權體制內，這只能把法院實際上置於立法機關之內，因此比 Kelsen\ix{persons}{kelsen@Kelsen, Hans（凱爾森）} 所建議的更強調憲法委員會的立法功能。

Ferejohn\ix{persons}{ferejohn@Ferejohn, John} 與 Pasquino\ix{persons}{pasquino@Pasquino, Pasquale} 列出該委員會審查權的範圍與限制：

\begin{enumerate}
\def\labelenumi{\arabic{enumi}.}
\tightlist
\item
  只能在制定法公布之前審查，且只在政府或（自 1974 年起）議員或參議員的重要政治少數提請時。
\item
  立法提案若被憲法委員會宣告違憲，便不能成為法律。
\item
  制定法一旦成為法律，憲法委員會無權審查；普通法院也無權在一般訴訟中進行憲法審查。
\item
  依法國的見解，制定法是主權者，不受任何外部檢驗，更不受法官或憲法法院檢驗；憲法法院作為立法機關自身的一部分，在保存法律主權的觀念上扮演關鍵角色。
\end{enumerate}

Ferejohn 與 Pasquino 這樣概括：「Legislative action cannot occur in the presence of an objection by the Conseil, but once it takes place, is unchecked by constitutional mechanisms」〔有憲法委員會的異議，立法行為不能成立；一旦成立，就不受憲法機制檢驗〕\pg{32 頁}。

因為法國的立法提案在憲法委員會拒絕時不能成為法律，審查（若被請求，且幾乎總是由政治少數針對政府所提立法提出）必須在立法行為之後立刻進行，面對的是其提案受到挑戰的現任政府。兩人認為，這表示審查必須迅速，而且鑑於法國政治制度與政治文化的多數決性質，處在可能高度政治化的情況\pg{32 頁}。Ferejohn 與 Pasquino 對照美國：憲法對制定法的挑戰須等到真正提出該議題的案件，往往在制定該法的立法機關消失之後很久；在德國、西班牙與義大利等後威權\ix{terms}{06@\zhlidxcat{審查制度之類型}!後威權政體}體制，憲法法院不能阻止法律生效，即使審查具有抽象性質，政府正在執行受爭議之法律，並沒有太大的政治壓力要迅速解決憲法問題（Ferejohn 與 Pasquino 並說政治壓力或許反向作用），抽象審查\ix{terms}{06@\zhlidxcat{審查制度之類型}!抽象審查與具體審查}因此可在政治戰火稍歇之後進行\pg{32 至 33 頁}。

關於組成，兩人認為或許因為位在立法機關之內，憲法委員會的組成相當獨特：成員只是偶爾有實質的法律或司法背景，法學教授也少；憲法委員會由任命權者的政治盟友主導。Ferejohn 與 Pasquino 補充這並不表示成員無法理解或提出法律論證，而是該機構看重的是政治判斷，這在他們必須迅速作成決定的急迫性之下尤為關鍵\pg{33 頁}。

\section{憲法法院的審議\pg{33 至 36 頁}}
\label{sec:FP-5}

\subsection{（一）三種制度與三種審議型態\pg{33 頁，註 18}}

第 4 節論證三種不同的憲法審查制度，產生三種可以區別的審議型態\pg{33 頁}。

\begin{longtable}[]{@{}
  >{\raggedright\arraybackslash}p{(\linewidth - 6\tabcolsep) * \real{0.1400}}
  >{\raggedright\arraybackslash}p{(\linewidth - 6\tabcolsep) * \real{0.4400}}
  >{\raggedright\arraybackslash}p{(\linewidth - 6\tabcolsep) * \real{0.3400}}
  >{\raggedright\arraybackslash}p{(\linewidth - 6\tabcolsep) * \real{0.0800}}@{}}
\caption{三種審查制度與其審議型態。}\tabularnewline
\toprule\noalign{}
\begin{minipage}[b]{\linewidth}\raggedright
制度
\end{minipage} & \begin{minipage}[b]{\linewidth}\raggedright
審查的所在與方式
\end{minipage} & \begin{minipage}[b]{\linewidth}\raggedright
審議型態
\end{minipage} & \begin{minipage}[b]{\linewidth}\raggedright
頁次
\end{minipage} \\
\midrule\noalign{}
\endhead
\bottomrule\noalign{}
\endlastfoot
權力分立\ix{terms}{06@\zhlidxcat{審查制度之類型}!權力分立體制}體制 & 憲法審查在普通法院進行；法官之間需要高度協調與溝通 & 傾向發展出公開的審議形式：公開聽審、公布意見與表決、先例理論 & 33 \\
Kelsen 型\ix{terms}{06@\zhlidxcat{審查制度之類型}!Kelsen 型模式}模式 & 憲法審查權集中於專門法庭；普通法院須能適用其裁判，故裁判須明確 & 傾向以單一聲音發言並闡明理由；很少公開聽審或公布不同意見\ix{terms}{04@\zhlidxcat{法院之功能}!不同意見與協同意見}書 & 33 \\
議會主權\ix{terms}{06@\zhlidxcat{審查制度之類型}!議會主權體制}體制 & 憲法審查是立法程序的一部分，目的在確保立法合於憲法要求；如何進行，取決於審查如何納入立法程序 & 依審查納入立法程序的方式而定 & 33 \\
\end{longtable}

Ferejohn\ix{persons}{ferejohn@Ferejohn, John} 與 Pasquino\ix{persons}{pasquino@Pasquino, Pasquale} 就議會主權體制中審查制度最為制度化的法國說明：憲法委員會對列入其審查事項的立法提案，在祕密會議中審議，沒有公開聽審或陳述\pg{33 頁}。法國憲法審查的重要特徵是，經議會多數通過的制定法在憲法委員會同意之前不能公布，所以審查程序承受巨大的時間壓力：延遲通常為一個月，「this period shall be reduced to eight days where the Government declares the matter urgent」〔政府宣告案件緊急時，此期間縮短為八日〕（註 18 引 1958 年 11 月第 58-1067 號條例第 25 條\ix{statutes}{03@法國!01@憲法委員會組織法（1958 年第 58-1067 號條例）!025@第 25 條}）。

\subsection{（二）法國憲法委員會的意見與解釋保留\pg{33 至 35 頁，註 19 至 22}}

憲法委員會以法院的名義發布意見，沒有紀錄表決，也沒有公布的不同意見\ix{terms}{04@\zhlidxcat{法院之功能}!不同意見與協同意見}書；意見通常非常簡短，不詳述對立法的憲法異議\pg{33 至 34 頁}。委員會可以報告提案全部或一部在憲法上有疑義，所依據的憲法淵源有數個：1789 年人權宣言、1946 年憲法前言，以及「Fundamental Principles underlying the laws of the {[}Third{]} Republic」〔第三共和國法律所依據的基本原則〕\pg{34 頁}。兩人稱最具影響力的許多裁判，涉及對行政機關過寬或不受限制的授權（註 20 舉車輛搜索案件與國有化裁判；1971 年發現結社自由是憲法原則的裁判也屬此型）。面對這類疑慮，政府通常可以重新起草立法，納入若干規範日後行政行為的標準\pg{34 至 35 頁}。

註 19 以 1971 年 7 月 16 日的意見為代表，說明該意見規定結社權並大幅修改委員會的角色，而且很短，並未真正闡明把 bloc de constitutionnalité\ix{terms}{06@\zhlidxcat{審查制度之類型}!憲法整體} 列為憲法審判依據的理由；註中並錄 J. Bell 的英譯（\en{French Constitutional Law}, Oxford: Clarendon Press, 1992, pp. 272--73），其中的理由段落稱：「among the fundamental principles recognized by the laws of the Republic and solemnly reaffirmed by the Constitution, is to be found the freedom of association」〔在共和國法律所承認並為憲法所莊嚴重申的基本原則中，可以找到結社自由〕\pg{34 頁}。

Ferejohn\ix{persons}{ferejohn@Ferejohn, John} 與 Pasquino\ix{persons}{pasquino@Pasquino, Pasquale} 接著列舉三點\pg{35 頁}：

\begin{enumerate}
\def\labelenumi{\arabic{enumi}.}
\tightlist
\item
  法國憲法委員會宣告立法提案無效的頻率不低，但更常認定提案的一部分與憲法原則相悖（註 21 引 A. Stone 的資料：1981 至 1993 年，約 50\% 被提請審查的法律被委員會指為有違憲之處）。
\item
  委員會愈來愈常發布意見，指出法律只有在以特定方式解釋時才在憲法上有效（註 22 引 Alexandre Viala, \en{Les réserves d'interprétation\ix{terms}{07@\zhlidxcat{憲法審判之審議}!解釋保留} dans la jurisprudence du Conseil constitutionnel}, Paris: L.G.D.J., 1999）。
\item
  Ferejohn 與 Pasquino 認為最後這項策略特別成問題，因憲法委員會與司法體系分離：「How can the Conseil offer any assurance that a promulgated law will be interpreted in a constitutional fashion?」〔憲法委員會如何能保證已公布的法律將以合憲的方式被解釋？〕沒有 stare decisis 之類的正式機制或法理可確保此事，唯一的保證是委員會向普通法院、最高法院（Cour de Cassation）或行政機關所提理由的說服力。兩人引述的跡象是：行政機關為尋求來自一個可能製造麻煩的法庭的明確指引，傾向把委員會的裁判用於制定對屬員的內部指引；普通法院為尋求法律穩定，常把委員會的解釋當作權威接受\pg{35 頁}。
\end{enumerate}

\subsection{（三）Kelsen 型與後威權法院的內部審議\pg{35 頁}}

Ferejohn\ix{persons}{ferejohn@Ferejohn, John} 與 Pasquino\ix{persons}{pasquino@Pasquino, Pasquale} 概括 Kelsen 型\ix{terms}{06@\zhlidxcat{審查制度之類型}!Kelsen 型模式}或後威權\ix{terms}{06@\zhlidxcat{審查制度之類型}!後威權政體}體制的法院：通常祕密審議，很少舉行公開聽審，意見以全體法院的名義發布，沒有記載表決，少有制度允許公布不同意見\ix{terms}{04@\zhlidxcat{法院之功能}!不同意見與協同意見}書\pg{35 頁}。直接開放訴訟當事人聲請的法院（如西班牙與德國）通常設有內部分工，以小組處理這類案件；有些法院准許或鼓勵就標的事項專業化（義大利常見，德國則反對）。Ferejohn 與 Pasquino 稱所有這些法院都是「internally deliberative」〔內部審議\ix{terms}{07@\zhlidxcat{憲法審判之審議}!內部審議}的〕，意即其工作多以合議與面對面的方式進行，真實地嘗試說服他人並作成集體說理的裁判；義大利尤其如此，法院以共識的裁判為目標而且通常達成（兩人註明這是從與法官及其他觀察者的討論所能得知的範圍）。

\subsection{（四）美國最高法院的外部審議\pg{35 至 36 頁}}

Ferejohn\ix{persons}{ferejohn@Ferejohn, John} 與 Pasquino\ix{persons}{pasquino@Pasquino, Pasquale} 寫道，美國最高法院的審議實踐與上述情形相當不同，可稱為 externally deliberative〔外部審議\ix{terms}{07@\zhlidxcat{憲法審判之審議}!外部審議}〕，與內部審議\ix{terms}{07@\zhlidxcat{憲法審判之審議}!內部審議}相對\pg{35 頁}。一部分原因在於法院對憲法解釋並無壟斷，主要作用在規範其他法院適用憲法的過程；這種協調或規制的角色，迫使法院以公開而透明的方式工作，使其他行動者，即法官、行政機關與律師，能清楚並（多半）成功地預期法院對尚未提到面前的案件會如何裁判。制定法可能在適用於個案時引起憲法問題，這使法院有時准許或鼓勵下級法院實驗，而不尋求立即解決限於文本本身的制定法的合憲性\pg{35 至 36 頁}。Ferejohn 與 Pasquino 說，這項策略只存在於准許具體審查\ix{terms}{06@\zhlidxcat{審查制度之類型}!抽象審查與具體審查}的體制，Kelsen 型\ix{terms}{06@\zhlidxcat{審查制度之類型}!Kelsen 型模式}體制僅在處理憲法訴願時可能發生（兩人列舉 amparo 與 Verfassungsbeschwerde）。

Ferejohn 與 Pasquino 列舉最高法院公開行事的幾個意義\pg{36 頁}：

\begin{enumerate}
\def\labelenumi{\arabic{enumi}.}
\tightlist
\item
  向一般訴訟當事人開放的案源，並公開其是否受理案件的決定；偶爾有成員甚至公開宣布個人對此的看法。
\item
  公開聽審是常態（至少對法院案件中最重要的部分），表決有紀錄，並公布不同意見\ix{terms}{04@\zhlidxcat{法院之功能}!不同意見與協同意見}書與協同意見書。
\item
  法院很少試圖以一個聲音發言，似乎寧可讓衝突與意見不一發酵（原文印作 ``ferment''）。
\end{enumerate}

就內部審議而言，Ferejohn 與 Pasquino 說「On most accounts the actual level of internal deliberation on the court may be low」〔依多數說法，法院內部審議的實際程度可能偏低〕：大法官在評議中不花多少時間審議案件中的具體爭點，觀察者通常報告大法官與自己的助理互動最多，與其他大法官較少，也很少以真正改變心意或票數為目的。內部審議多半發生於評議時的初步表決之後，即大法官決定能否全部或部分同意為法院所撰意見時；這種審議大多以書面進行，並非面對面的互動。兩人認為正因如此，審議過程對局外人比面對面的審議透明得多：「The public manner of conducting its business puts the Court at the center of wider deliberative processes within the judiciary, the legal, and political communities」〔公開辦理事務的方式，使法院處於司法、法律與政治各社群中更廣的審議過程的中心〕\pg{36 頁}。

圖~\ref{fig:FP-deliberation} 依三種審查制度，並列審議在法院之內進行或對著法院之外的行動者進行。

\begin{qafigure}[審議的所在與對象]{審議的所在與對象。Ferejohn 與 Pasquino，33 至 36 頁。\label{fig:FP-deliberation}}
% Ferejohn 與 Pasquino 第 4 節（33 至 36 頁）：審議發生在法院之內或對著法院之外。
% Kelsen 型與後威權法院：祕密審議、意見以全體法院名義發布（35 頁）；法官面對面說服、產生集體說理的裁判（35 頁）；
%   裁判須明確，普通法院始能適用（33 頁）；很少公開聽審或公布不同意見書（33 頁）。
% 法國憲法委員會：祕密會議、無公開聽審（33 頁）；意見簡短、無表決紀錄（33 至 34 頁）；與司法體系分離，
%   無 stare decisis 之類機制，唯一的保證是向普通法院、Cour de Cassation、行政機關所提理由的說服力（35 頁）。
% 美國最高法院：公開聽審、記名表決、公布協同與不同意見書（36 頁）；內部審議多在初步表決後以書面進行（36 頁）；
%   公開透明使法官、行政機關、律師得預期法院之裁判（35 頁）；法院處於司法、法律、政治社群更廣的審議過程中心（36 頁）。
% 色調：外部審議之法院與其流向用 accent（terracotta），其餘 slate；卡片淡底加左側色條。
\definecolor{figRev}{HTML}{B08060}\definecolor{figRevbg}{HTML}{F5ECEB}
\definecolor{figLine}{HTML}{E8D3D1}
\providecommand{\placecard}[6]{%
  \fill[#1bg] (#2,#3) rectangle (#4,#5);
  \fill[#1] (#2,#3) rectangle ($(#2,#5)+(0.7,0)$);
  \node[t] at ($(#2,#3)+(2.4,-1.9)$) {#6};}
\begin{tikzpicture}[qa,x=1mm,y=1mm,
    t/.style={anchor=north west,inner sep=0pt,align=left,execute at begin node={\raggedright\setstretch{1.2}}},
    note/.style={font=\scriptsize,text=dimink,inner sep=0pt,align=left,anchor=west,execute at begin node={\setstretch{1.15}}},
    link/.style={figGray,line width=0.6pt,-{Stealth[length=4pt,width=3.2pt]}},
    weak/.style={figGray,line width=0.6pt,dash pattern=on 2pt off 1.6pt,-{Stealth[length=4pt,width=3.2pt]}},
    flow/.style={figRev,line width=1.1pt,-{Stealth[length=4.5pt,width=3.6pt]}},
    both/.style={figGray,line width=0.5pt,{Stealth[length=3pt,width=2.6pt]}-{Stealth[length=3pt,width=2.6pt]}}]
  \foreach \x in {47.5,98.5}{\draw[figLine,line width=0.4pt] (\x,5) -- (\x,-96);}
  % 欄頭
  \foreach \x/\h/\s in {
      0/{Kelsen 型與後威權法院}/{內部審議（35 頁）},
      51/法國憲法委員會/{議會主權體制（33 頁）},
      102/美國最高法院/{外部審議（35 頁）}}{
    \node[t,font=\footnotesize] at (\x,4) {\h};
    \node[t,font=\scriptsize,text=dimink] at (\x,-0.6) {\s};
  }
  % Kelsen 型：審議在法官之間
  \begin{scope}[shift={(0,0)}]
    \placecard{figGray}{0}{-12}{44}{-46}{憲法法院\\{\scriptsize 祕密審議（35 頁）}\\{\scriptsize 很少公開聽審（33 頁）}\\{\scriptsize 意見以全院名義發布（35 頁）}}
    \foreach \p in {10,22,34}{\fill[figGray] (\p,-35) circle (1.3);}
    \draw[both] (11.8,-35) -- (20.2,-35);
    \draw[both] (23.8,-35) -- (32.2,-35);
    \node[note,anchor=north west,text=caseink] at (2.4,-37.4) {面對面說服，集體說理（35 頁）};
    \draw[link] (8,-46) -- (8,-68);
    \node[note] at (11,-57) {裁判須明確，普通法院\\始能適用（33 頁）};
    \placecard{figGray}{0}{-68}{20}{-80}{{\scriptsize 普通法院}}
  \end{scope}
  % 法國：祕密審議，意見向外傳遞但無正式拘束機制
  \begin{scope}[shift={(51,0)}]
    \placecard{figGray}{0}{-12}{44}{-46}{憲法委員會\\{\scriptsize 祕密會議，無公開聽審（33 頁）}\\{\scriptsize 意見簡短（33 至 34 頁）}\\{\scriptsize 與司法體系分離（35 頁）}}
    \draw[figGray,line width=0.6pt,dash pattern=on 2pt off 1.6pt] (7,-46) -- (7,-59);
    \draw[figGray,line width=0.6pt,dash pattern=on 2pt off 1.6pt] (7,-59) -- (37,-59);
    \draw[weak] (7,-59) -- (7,-68);
    \draw[weak] (22,-59) -- (22,-68);
    \draw[weak] (37,-59) -- (37,-68);
    \node[note] at (10,-52.5) {無正式機制，唯憑理由\\的說服力（35 頁）};
    \placecard{figGray}{0}{-68}{14}{-80}{{\scriptsize 普通法院}}
    \placecard{figGray}{15}{-68}{29}{-80}{{\scriptsize Cour de}\\{\scriptsize Cassation}}
    \placecard{figGray}{30}{-68}{44}{-80}{{\scriptsize 行政機關}}
  \end{scope}
  % 美國：公開行事，對著法院之外的行動者與更廣的社群
  \begin{scope}[shift={(102,0)}]
    \placecard{figRev}{0}{-12}{44}{-46}{最高法院\\{\scriptsize 公開聽審，記名表決（36 頁）}\\{\scriptsize 公布協同與不同意見書（36 頁）}\\{\scriptsize 內部審議多在初步表決後}\\{\scriptsize 以書面進行（36 頁）}}
    \draw[flow] (7,-46) -- (7,-59);
    \draw[figRev,line width=1.1pt] (7,-59) -- (37,-59);
    \draw[flow] (7,-59) -- (7,-68);
    \draw[flow] (22,-59) -- (22,-68);
    \draw[flow] (37,-59) -- (37,-68);
    \node[note] at (10,-52.5) {公開行事，使其得預期\\裁判（35 頁）};
    \placecard{figGray}{0}{-68}{14}{-80}{{\scriptsize 下級法院}}
    \placecard{figGray}{15}{-68}{29}{-80}{{\scriptsize 行政機關}}
    \placecard{figGray}{30}{-68}{44}{-80}{{\scriptsize 律師}}
    \placecard{figGray}{0}{-84}{44}{-96}{{\scriptsize 法院處於司法、法律、政治社群}\\{\scriptsize 更廣的審議過程之中（36 頁）}}
  \end{scope}
\end{tikzpicture}

\end{qafigure}
